\documentclass{training}
% Filigrane PAR MODULE : \uxfond (PNG A4 dilue 30 %). Cadre noir vectoriel commun.
% Proportions original 744 px : marge 20/744, trait 1/744. Exterieur BLANC.
% Cadre via eso-pic +\put (coordonnees absolues) — pas de remember picture (decalage).
\usepackage{eso-pic}

\providecommand{\uxfond}{../../resources/fond/norme_fond_paliers_30.png}

\newlength{\uxmargin}
\newlength{\uxthick}
\newlength{\uxinnerw}
\newlength{\uxinnerh}
\setlength{\uxmargin}{0.02688\paperwidth}% 20/744
\setlength{\uxthick}{0.001344\paperwidth}% 1/744
\setlength{\uxinnerw}{\dimexpr\paperwidth-2\uxmargin-2\uxthick\relax}
\setlength{\uxinnerh}{\dimexpr\paperheight-2\uxmargin-2\uxthick\relax}

% Contenu strictement a l'interieur du trait (+ jour confortable).
\geometry{%
  includeheadfoot,
  left=22mm,
  right=22mm,
  top=20mm,
  bottom=18mm,
  headheight=30pt,
  headsep=5mm,
  footskip=12mm%
}

\AddToShipoutPictureBG{%
  \AtPageLowerLeft{%
    % Fond blanc page entiere
    \put(0,0){\color{white}\rule{\paperwidth}{\paperheight}}%
    % Filigrane module a l'interieur du cadre
    \put(\LenToUnit{\dimexpr\uxmargin+\uxthick\relax},\LenToUnit{\dimexpr\uxmargin+\uxthick\relax}){%
      \includegraphics[width=\uxinnerw,height=\uxinnerh]{\uxfond}}%
    % Cadre noir : bas, haut, gauche, droite
    \put(\LenToUnit{\uxmargin},\LenToUnit{\uxmargin}){%
      \color{black}\rule{\dimexpr\paperwidth-2\uxmargin\relax}{\uxthick}}%
    \put(\LenToUnit{\uxmargin},\LenToUnit{\dimexpr\paperheight-\uxmargin-\uxthick\relax}){%
      \color{black}\rule{\dimexpr\paperwidth-2\uxmargin\relax}{\uxthick}}%
    \put(\LenToUnit{\uxmargin},\LenToUnit{\uxmargin}){%
      \color{black}\rule{\uxthick}{\dimexpr\paperheight-2\uxmargin\relax}}%
    \put(\LenToUnit{\dimexpr\paperwidth-\uxmargin-\uxthick\relax},\LenToUnit{\uxmargin}){%
      \color{black}\rule{\uxthick}{\dimexpr\paperheight-2\uxmargin\relax}}%
  }%
}

\makeatletter
\fancyhf{}
\fancyheadoffset{0pt}
\renewcommand{\headrulewidth}{0pt}
\lhead{\parbox[t]{0.56\textwidth}{\RaggedRight\footnotesize\sloppy\@title}}
\rhead{\parbox[t]{0.36\textwidth}{\RaggedLeft\footnotesize\sloppy\@subtitle}}
\cfoot{\thepage}
% Coupures : underscore et slash dans les noms de fichiers / consignes
\DeclareRobustCommand{\_}{\textunderscore\allowbreak}
\makeatother

% Evite les depassements de longues lignes (Github, Name_FirstName_...)
\sloppy
\emergencystretch=3em
\hbadness=10000

% Sommaire type 42 : chapitres seulement
\setcounter{tocdepth}{0}
\renewcommand{\cftchappresnum}{}
\renewcommand{\cftchapaftersnum}{}
\setlength{\cftchapnumwidth}{2.0em}
\setlength{\cftbeforechapskip}{0.35em}
\cftsetindents{chapter}{0em}{2.0em}


\begin{document}

\schoollogo{../../resources/jeunecrack.png}
\title{Building an interface}
\subtitle{Piscine UX/UI --- UI}
\member{Piscine UX/UI}{https://github.com/01-edu/UX-UI}
\summary{A styleguide, a library, screens, motion, and a desirability test.}
\version{1.0}
\prerequisite{The previous module of this piscine, or none if this is your first.}
\maketitle

\newpage
\tableofcontents
\newpage


\chapter{Instructions}

The foreword is not an exercise. It will not be graded. The exercises start at the next chapter, from the easiest contract to the hardest. A later exercise is not taken into account if an earlier one does not hold.

Only what is in the turn-in directory of an exercise is looked at. Do not add files that the exercise does not name. The name of a deliverable follows the pattern in the exercise: \texttt{Name\_FirstName\_DeliverableName\_Date\_VersionNumber}.

You will not be given the gesture. The examples in an exercise sometimes carry a constraint that the sentence does not repeat. Read them. Your other references are the manual of the tool you picked, the person sitting next to you, and the links printed in the exercise.

The order is the order of the quest. Skipping a step because a later screen looks more interesting produces a screen nobody can explain.

\subsection*{About the pictures}

The mark on the cover is the project logo. The picture in the foreword is the illustration of this module. The pale ground behind the text is the page watermark. None of them is a hint.

\subsection*{About machines that write}

You may use a tool to look up a term, the way you would use a dictionary. You may not replace an interview, a test, or a file with a generated text and sign it. If you cannot point at the person, the screen, or the line you claim, the exercise is not done.


\chapter{Foreword}

Gary Hustwit's \emph{Objectified} (2009) films people who argue about a handle, a toothpick, a chair. The argument is the work. This quest is that argument, held on screens you draw yourself. There is a second shelf in the foreword, an AI playlist. It is there so you know what the tools sound like. It is not a licence to submit a generated deck.

\begin{center}
\includegraphics[width=0.78\textwidth,height=0.32\textheight,keepaspectratio]{../../resources/intros/UI-03-building-an-interface.jpg}\\[0.4em]
{\footnotesize Picture: freddie marriage fredmarriage / CC0}
\end{center}

Recommended watching. The videos are context. They are not the deliverable.
\begin{itemize}
\item UI - Introduction \& Quest n°1, 01Edu \url{https://www.youtube.com/watch?v=UbE6WMeuvQg&list=PLHyAJ_GrRtf-_aBft2jQL4-QAunn2bLvu}
\item Claude AI Enablement piscine, 01Edu \url{https://www.youtube.com/watch?v=M94thLMc6vA&list=PLCiJelRloNB8}
\end{itemize}
The last playlist is an AI piscine (\emph{Claude AI Enablement}). Watch it if you want to hear how a model is sold as a design assistant. The files you turn in still have to come from your research, your screens, and your tests. A generated board that you cannot walk through does not count.

\chapter{Exercise 00: UI IV - Ex 1 - Styleguide}
\exnumber{00}
\extitle{UI IV - Ex 1 - Styleguide}
\exfiles{ex00/Figma file, named Name\_FirstName\_DeliverableName\_Date\_VersionNumber}
\exauthorize{Those named in the exercise}
\makeheaderfiles

\textbf{Context:}

In this quest, you'll have to build your own interface.

After designing your own style guide with a moodboard, colors and typography, you can use the series of consistency principles, atomic design structure, the techniques to build you own library and then a series of screens. Step by step, you'll have built an animated high-fidelity prototype.

The very last step will be to test it on users to detect the impressions they get when they discover your app.

Once again, feel free to get as much inspiration as you need from existing design systems. Dribbble, consistency and structure are your best friends!

\textbf{Instructions:}

Your agency works on a digital product to create an e-commerce app for second-hand electronic devices. You collaborate with a UX designer from your team. They have deisgned the wireframes and pass it on to you!

Design a style guide (= a moodboard with typography + examples of typography with buttons, title, description and content).

Also think about the values and spirit you want to spread through this moodboard, we'll get back to it in Ex. 5.

\textbf{Deliverables:}

\begin{itemize}
\item A Figma moodboard that includes:
\begin{itemize}
\item Pictures
\item 2 to 3 main colors
\item Examples of typography applied to graphic elements
\begin{itemize}
\item Title
\item Description or subtitle
\item Content
\item Button
\end{itemize}
\end{itemize}
\item Upload the Figma file on Github, labeled as follows: "Name\_FirstName\_DeliverableName\_Date\_VersionNumber". For example, it can be labeled as "Doe\_John\_MidFiPrototype\_05242024\_V1".
\item Add a title within the file.
\item Pay attention to the global aspect of the file. It must be clear, simple and easy to read. You can get inspiration on canva.com but don't overload your design with too much details!
\end{itemize}

\textbf{Resources:}

\begin{itemize}
\item Video tutorial on Figma \url{https://www.youtube.com/watch?v=FTFaQWZBqQ8}
\item Beginning graphic design \url{https://edu.gcfglobal.org/en/beginning-graphic-design/}
\item Gareth graphic studio \url{https://www.youtube.com/c/GarethDavidStudio}
\item Example of a Style Guide in UI \url{https://www.pinterest.pt/pin/454019206175050454/}
\end{itemize}


\chapter{Exercise 01: UI IV - Ex 2 - Design system library}
\exnumber{01}
\extitle{UI IV - Ex 2 - Design system library}
\exfiles{ex01/Figma file, named Name\_FirstName\_DeliverableName\_Date\_VersionNumber}
\exauthorize{Those named in the exercise}
\makeheaderfiles

\textbf{Instructions:}

Based on the moodboard you designed in Ex 1, create a design system with at least 30 elements, including atoms, molecules and organisms, and include multistates.

\textbf{Deliverables:}

\begin{itemize}
\item A Figma file with a library including:
\begin{itemize}
\item At least graphic design 30 elements
\item At least 5 atoms
\item At least 5 molecules
\item At least 5 organisms
\item At least 10 multistate variations
\end{itemize}
\item Upload the Figma file on Github, labeled as follows: "Name\_FirstName\_DeliverableName\_Date\_VersionNumber". For example, it can be labeled as "Doe\_John\_MidFiPrototype\_05242024\_V1".
\item Add a title within the file.
\item Pay attention to the global aspect of the file. It must be clear, simple and easy to read. You can get inspiration on canva.com but don't overload your design with too much details!
\end{itemize}

\textbf{Tips:}

\begin{itemize}
\item Start by asking yourself what key elements may appear in an e-commerce app.
\item Feel free to get inspiration from existing design systems. Browse Dribble and the Figma community!
\end{itemize}

\textbf{Resources:}

\begin{itemize}
\item FIGMA FILE (to duplicate) \url{https://www.figma.com/file/U26mctkrOqxiUeAlNHzSxW/UI-IV---Ex-3---E-COMMERCE%3A-Design-screens!?node-id=0%3A1}
\item Introduction to design systems \url{https://www.youtube.com/watch?v=xEuBCUngJ_U}
\item Create a design system - Figma course \url{https://www.youtube.com/watch?v=RYDiDpW2VkM}
\item Building a design system library - Lyft \url{https://medium.com/tap-to-dismiss/building-a-design-system-library-532ef2492811}
\item Examples of UI Buttons library \url{https://www.pinterest.co.kr/pin/663014376380850291/}
\end{itemize}


\chapter{Exercise 02: UI IV - Ex 3 - Design screens}
\exnumber{02}
\extitle{UI IV - Ex 3 - Design screens}
\exfiles{ex02/Figma file, named Name\_FirstName\_DeliverableName\_Date\_VersionNumber}
\exauthorize{Those named in the exercise}
\makeheaderfiles

\textbf{Instructions:}

Based on the wireframes and your design system library, take your mid-fidelity prototype to a high-fidelity level!

Select 5 screens that will show the user flow of a customer who:

\begin{itemize}
\item arrives on the homepage
\item browses the content
\item checks one product page
\item adds the product to the cart
\item checks out
\end{itemize}

Use your components, call to actions and visuals to complete the wireframes.

\textbf{Deliverables:}

\begin{itemize}
\item 5 high-fidelity screens from Figma
\item Upload the Figma file on Github, labeled as follows: "Name\_FirstName\_DeliverableName\_Date\_VersionNumber". For example, it can be labeled as "Doe\_John\_MidFiPrototype\_05242024\_V1".
\item Add a title within the file.
\item Pay attention to the global aspect of the file. It must be clear, simple and easy to read. You can get inspiration on canva.com but don't overload your design with too much details!
\end{itemize}

\textbf{Resources:}

\begin{itemize}
\item Unsplash \url{https://unsplash.com/} for visuals
\item From low-fidelity to high-fidelity \url{https://mockitt.wondershare.com/prototyping/low-fidelity-vs-high-fidelity.html}
\end{itemize}


\chapter{Exercise 03: UI IV - Ex 4 - Animation}
\exnumber{03}
\extitle{UI IV - Ex 4 - Animation}
\exfiles{ex03/Figma file, named Name\_FirstName\_DeliverableName\_Date\_VersionNumber}
\exauthorize{Those named in the exercise}
\makeheaderfiles

\textbf{Instructions:}

Animate your 5 screens on Figma. Think about multistate!!

\textbf{Deliverable:}

\begin{itemize}
\item An animated high-fidelity prototype on Figma
\item Upload the Figma file on Github, labeled as follows: "Name\_FirstName\_DeliverableName\_Date\_VersionNumber". For example, it can be labeled as "Doe\_John\_MidFiPrototype\_05242024\_V1".
\item Add a title within the file.
\item Pay attention to the global aspect of the file. It must be clear, simple and easy to read. You can get inspiration on canva.com but don't overload your design with too much details!
\end{itemize}

\textbf{Resources:}

\begin{itemize}
\item Efficiently Manage Your Designs - How Mockplus State Page Helps To View \& Manage Multiple State \url{https://help.mockplus.com/p/372}
\end{itemize}

\textbf{Quote:}

'Speak human: Use the words your users use.' Dan Saffer | Microinteractions


\chapter{Exercise 04: UI IV - Ex 5 - Desirability testing}
\exnumber{04}
\extitle{UI IV - Ex 5 - Desirability testing}
\exfiles{ex04/zip archive, named Name\_FirstName\_DeliverableName\_Date\_VersionNumber}
\exauthorize{Those named in the exercise}
\makeheaderfiles

\textbf{Instructions:}

Run a desirability test!

The point of this phase is to make sure the values and the spirit you chose for your design.

\begin{itemize}
\item Get back to the values and spirit you chose in Ex 1. Extract 2 to 5 adjectives.
\item Prepare a board with 25+ adjectives, randomly distributed.
\item Show your design to 5 different people, and ask them to select the adjectives among the list from the board that they are inspired with.
\item If the selected adjectives chosen by the users are consistent with the ones you had chosen, congratulations!
\item If the selected adjectives chosen by the users are not consistent, that's fine! Make 3 suggestions about how to make your moodboard and library evolve.
\end{itemize}

\textbf{Deliverables:}

\begin{itemize}
\item A desirability testing board with 25+ different adjectives
\item Pictures from the desirability test
\item A PDF document with the outcomes of the desirability test
\item If the outcomes are not consistent with the values, 3 suggestions on how to make the moodboard and design library evolve.
\item Upload the documents on Github, in a zip folder named "ProjectTitle\_Name\_FirstName", with all project deliverables as follows: "Name\_FirstName\_DeliverableName\_Date\_VersionNumber".  For example, the first deliverable can be named as "Doe\_John\_MidFiPrototype\_05242024\_V1".
\item Add a title within all written documents.
\item Pay attention to the global aspect of the deliverables. They must be clear, simple and easy to read. You can get inspiration on canva.com but don't overload your design with too much details!
\end{itemize}

\textbf{Tips:}

\begin{itemize}
\item Avoid confirmation biais: don't tell them the values and ask them to approve them!
\item Good to know: Desirability testing is great to assess the impressions a design gives to the users. Whereas Usability testing is great to understand what users understand your prototype and how they navigate across the product.
\end{itemize}

\textbf{Resources:}

\begin{itemize}
\item Desirability Testing analyzing an emotional response to a design \url{https://www.slideshare.net/megangrocki/desirability-testing-analyzing-emotional-response-to-a-design-11262575}
\end{itemize}

\textbf{Quote:}

\begin{itemize}
\item 'Interaction design focuses on the design of behavior of a person's interaction with a digital interface. It is also concerned with satisfying the needs and desires of the people who will interact with a product or service.' Alan Cooper, About Face: The Essentials of Interaction Design
\end{itemize}


\end{document}
